\section{Discussion}
\label{sec:discussion}

\subsection{Why INT8 Does Not Degrade Retrieval Quality}
\label{sec:rkdef}

The zero-degradation result ($\Delta$R@5\,=\,0, $\Delta$MRR\,=\,0 on
Spider~1.0 across 1{,}034 queries; formal definitions in
Section~\ref{sec:metrics}) admits a straightforward empirical explanation.
Dynamic INT8 quantization replaces each FP32 weight with its nearest
8-bit integer under a per-tensor scale $s = \max(|W|)/127$, introducing
rounding errors bounded by $s/2$ per element.
For arctic-embed-m's 768-dimensional embeddings trained with contrastive
objectives, semantic similarity is distributed across the entire embedding
dimension: no single dimension carries a disproportionate share of the
discriminative signal.
Consequently, the per-element rounding noise averages out in the
inner-product similarity computation, and the relative ordering of
table candidates is preserved.
This result is consistent with prior work on INT8 quantization of BERT
encoders~\citep{zafrir2019q8bert}, which reported comparable stability on
GLUE tasks, and generalises here to the cross-domain schema retrieval
setting.

This observation generalizes: embedding models trained for retrieval
with large hidden dimensions are robust to INT8 quantization.
Practitioners deploying dense retrievers on edge hardware
(Jetson Nano, Raspberry~Pi~4, mobile devices) should consider INT8
quantization as a default optimization, as it uniformly reduces RAM
and improves CPU latency without retrieval-quality risk.

\subsection{Synergy Between Self-Consistency and Two-Pass Correction}

Self-consistency and two-pass correction address complementary failure
modes.
Self-consistency reduces \textit{systematic} errors: when three
independent candidates are generated, at least one is likely to avoid
a particular structural mistake.
Two-pass correction reduces \textit{runtime} errors: queries that parse
successfully but fail at execution (wrong column name, type mismatch)
are correctable given the error message.
The synergy observed in Table~\ref{tab:ablation} (+3.87\,pp combined
vs.\ +2.45\,pp and +3.71\,pp individually) confirms that these
two components target distinct error types.

\subsection{Spider~2.0-Lite Difficulty}

The low absolute accuracy on Spider~2.0-Lite (EX\,$\approx$\,0.10)
reflects structural differences from Spider~1.0.
Spider~2.0-Lite queries require window functions, recursive CTEs,
and five-way analytical joins; these constructs appear rarely in the
24-example gold SQL pool, limiting the effectiveness of few-shot retrieval.
The 24-example pool contains only 16 unique databases, leaving 34\%
of test questions without a same-database few-shot example.

A larger pool, task-specific fine-tuning of the LLM, or schema-aware
SQL decomposition represent the most promising directions
for improving Spider~2.0-Lite accuracy beyond the current results.

\subsection{Threats to Validity}
\label{sec:threats}

Several threats to validity apply to the experimental conclusions.

\textbf{Platform dependence.}
INT8 latency results are specific to the qnnpack backend on ARM64 (Apple
M-series).
Performance on x86-64 (fbgemm backend) or CUDA may differ; in particular,
NVIDIA Tensor Cores provide native INT8 SIMD at higher throughput,
making the latency gap between INT8 and FP32 on GPU hardware narrower
than reported here.

\textbf{LLM non-determinism.}
SQL generation uses self-consistency at temperature $T\,=\,0.7$, introducing
stochastic variation across runs.
The 95\% bootstrap confidence intervals (Table~\ref{tab:ablation}) account
for this at the population level, but individual run-to-run variance
may differ from reported means by up to $\pm$0.3\,pp.

\textbf{Downstream EX comparison.}
The FP32 best-configuration run (EX\,=\,0.824) and the INT8 retriever
run were executed on different hardware (local Mac vs.\ server GPU), with
the same Ollama LLM.
A fully controlled comparison---identical hardware, identical LLM version,
identical random seed---is left for future work.

\textbf{Spider~2.0-Lite coverage.}
Only 135 of 547 Spider~2.0-Lite questions use local SQLite databases;
the remaining 412 require BigQuery or Snowflake access.
Results on the local subset may not generalize to the full benchmark,
particularly for cloud-specific SQL dialects.

\textbf{Few-shot pool size.}
The Spider~2.0-Lite few-shot pool contains 24 gold SQL examples across
16 databases; 34\% of the 135 test questions have no pool example from
the same database.
R@K estimates on Spider~2.0-Lite carry uncertainty due to this small
sample size and should be interpreted as indicative rather than definitive.

\subsection{Limitations}

Four limitations bound the generalizability of these findings.
First, Qwen3.5~9B generates SQL that contains semantic errors on complex
queries; replacing it with a larger or fine-tuned code model is the most
direct path to closing the remaining gap to DAIL-SQL on Spider~1.0 and
improving Spider~2.0-Lite accuracy beyond 0.111.
Second, the 24-example few-shot pool for Spider~2.0-Lite is too small to
provide diverse coverage across 30 databases, limiting the benefit of
few-shot retrieval on this benchmark.
Third, a fully controlled downstream EX comparison between FP32 and INT8
retrieval requires both configurations to run on identical hardware and
LLM versions; the retrieval-quality parity ($\Delta$R@5\,=\,0) provides
strong indirect evidence of $\Delta$EX\,=\,0, but direct measurement is
left for future work.
Fourth, Spider~2.0-Lite results cover only the 135-question local SQLite
subset; the full 547-question benchmark requires BigQuery and Snowflake
access, and the SQL dialects used there may respond differently to the
proposed pipeline.
